\chapter{Triglycerides}
\index{Triglycerides}

\section*{Definition and Overview}
Triglycerides, also known as triacylglycerols (TAG), are a class of lipids that represent the most abundant form of fat in the human body. Chemically, a triglyceride molecule consists of a single glycerol backbone esterified to three fatty acid chains. These fatty acids can vary in chain length and degree of saturation (ranging from saturated to monounsaturated and polyunsaturated), which in turn determines the physical properties of the lipid, such as its melting point. In human physiology, triglycerides serve as the primary energy reservoir. While glycogen stored in the liver and skeletal muscles provides a rapidly accessible but limited supply of carbohydrates, triglycerides stored in adipose tissue represent a highly concentrated, long-term energy reserve capable of sustaining the body during prolonged periods of fasting or increased metabolic demand.

Because triglycerides are hydrophobic and insoluble in water, they cannot circulate freely in the blood plasma. Instead, they must be packaged and transported within specialized spherical macromolecules known as lipoproteins. These particles possess a hydrophobic core composed of triglycerides and cholesteryl esters, surrounded by a hydrophilic monolayer of phospholipids, free cholesterol, and apolipoproteins. The metabolic pathway of triglycerides is divided into two primary systems: the exogenous pathway, which transports dietary fats absorbed by the small intestine in the form of large, triglyceride-rich chylomicrons, and the endogenous pathway, which distributes lipids synthesized by the liver in the form of very-low-density lipoproteins (VLDL). As these particles circulate, the enzyme lipoprotein lipase (LPL), located on the capillary endothelial walls of adipose and muscle tissues, hydrolyzes the triglycerides, releasing free fatty acids for cellular energy production or storage, and leaving behind denser, cholesterol-rich lipoprotein remnants.

In clinical practice, monitoring serum triglyceride levels is a fundamental component of assessing cardiovascular and metabolic health. Hypertriglyceridemia, defined as an abnormal elevation of fasting plasma triglyceride levels, is a common dyslipidemia. While historical debates questioned whether triglycerides were directly atherogenic, modern genetic, epidemiological, and clinical trial data have established that elevated levels of triglyceride-rich lipoproteins and their remnants are independent risk factors for atherosclerotic cardiovascular disease (ASCVD). Furthermore, severe elevations in triglycerides pose a direct, acute risk for the development of acute pancreatitis, a potentially life-threatening inflammatory condition of the pancreas.

\section*{Epidemiology}
Hypertriglyceridemia is highly prevalent globally, particularly in industrialized nations where sedentary lifestyles and hypercaloric diets are common. According to data from the National Health and Nutrition Examination Survey (NHANES) in the United States, approximately 25\% to 30\% of the adult population has triglyceride levels above the normal threshold of 150 mg/dL (1.7 mmol/L). The prevalence of this condition is closely linked to the rising rates of obesity, metabolic syndrome, and type 2 diabetes.

The epidemiological distribution of triglyceride levels varies significantly across demographics:
\begin{itemize}
    \item \textbf{Age}: Average serum triglyceride levels tend to increase with age in both men and women. This increase is typically associated with age-related weight gain, increases in visceral adiposity, decreased physical activity, and a gradual decline in the metabolic clearance of lipoproteins.
    \item \textbf{Biological Sex}: In younger and middle-aged populations, men generally have higher triglyceride levels than premenopausal women. This disparity is primarily due to the metabolic effects of endogenous estrogens, which enhance the clearance of VLDL particles and modulate hepatic lipid production. Following menopause, however, the decline in estrogen levels is associated with a significant increase in triglyceride levels in women, often matching or exceeding the prevalence seen in men of similar age.
    \item \textbf{Ethnicity}: Epidemiological studies have revealed distinct differences in lipid profiles among ethnic groups. For instance, Hispanic populations (particularly those of Mexican descent) demonstrate a high prevalence of hypertriglyceridemia, which correlates with higher rates of insulin resistance, abdominal obesity, and metabolic dysfunction-associated steatotic liver disease (MASLD). Conversely, non-Hispanic Black individuals consistently show lower average triglyceride levels than their white counterparts, even when adjusted for body mass index. This is partly explained by higher levels of lipoprotein lipase activity and lower hepatic secretion of VLDL in Black populations.
\end{itemize}

\section*{Aetiology and Pathophysiology}
The underlying causes of hypertriglyceridemia can be categorized into primary (genetic) and secondary (acquired or environmental) factors. The pathophysiology is centered on either the overproduction of triglyceride-rich VLDL by the liver, the impaired clearance of chylomicrons and VLDL from the circulation, or a combination of both.

Primary genetic causes represent mutations in key proteins involved in lipid metabolism:
\begin{itemize}
    \item \textbf{Familial Chylomicronemia Syndrome (FCS)}: A rare, autosomal recessive disorder characterized by a severe, complete deficiency of lipoprotein lipase (LPL) or one of its regulatory cofactors, such as apolipoprotein C-II (ApoC-II), apolipoprotein A-V (ApoA-V), glycosylphosphatidylinositol-anchored lipoprotein-binding protein 1 (GPIHBP1), or lipase maturation factor 1 (LMF1). In these patients, the body cannot clear chylomicrons from the blood, resulting in extreme hypertriglyceridemia (often exceeding 1000--2000 mg/dL) from infancy or childhood.
    \item \textbf{Familial Combined Hyperlipidemia (FCHL)}: A common, polygenic disorder characterized by elevated levels of both cholesterol and triglycerides, accompanied by an increase in apolipoprotein B (ApoB). The primary defect involves the hepatic overproduction of VLDL, leading to an increased number of circulating atherogenic particles and a high risk of premature coronary artery disease.
    \textbf{Familial Hypertriglyceridemia (FHTG)}: An autosomal dominant condition characterized by moderate elevations in triglycerides (typically 200--500 mg/dL) with normal cholesterol and ApoB levels. The pathophysiology involves the production of large, triglyceride-rich VLDL particles with normal particle numbers.
    \item \textbf{Dysbetalipoproteinemia (Type III Hyperlipoproteinemia)}: An autosomal recessive disorder typically associated with homozygosity for the apolipoprotein E2 (APOE2/E2) isoform. Because the ApoE2 isoform has a low affinity for LDL receptors, the clearance of chylomicron remnants and VLDL remnants (intermediate-density lipoproteins) is severely impaired, causing significant elevations in both cholesterol and triglycerides.
\end{itemize}

Secondary acquired causes are far more common and represent the vast majority of clinical cases:
\begin{itemize}
    \item \textbf{Dietary and Lifestyle Factors}: High intake of simple carbohydrates, refined sugars (especially fructose), and saturated fats stimulates hepatic de novo lipogenesis, raising VLDL production. Alcohol consumption is another potent trigger; ethanol metabolism alters the hepatic redox state (increasing the NADH/NAD+ ratio), which suppresses fatty acid oxidation and promotes triglyceride synthesis.
    \item \textbf{Obesity and Insulin Resistance}: In insulin-resistant states, the normal inhibitory effect of insulin on hormone-sensitive lipase in adipose tissue is impaired. This leads to an excessive release of free fatty acids into the portal circulation, which are taken up by the liver, re-esterified into triglycerides, and secreted as VLDL. Insulin resistance also reduces the tissue expression and activity of LPL, further slowing clearance.
    \item \textbf{Diabetes Mellitus}: In poorly controlled type 1 or type 2 diabetes, insulin deficiency or resistance impairs the activation of LPL, leading to decreased clearance of chylomicrons and VLDL.
    \item \textbf{Hypothyroidism}: Thyroid hormones stimulate the expression of LDL receptors and LPL. Hypothyroidism reduces the clearance of VLDL remnants, elevating both triglycerides and LDL cholesterol.
    \item \textbf{Renal Diseases}: In nephrotic syndrome, the liver increases lipogenesis in response to hypoalbuminemia and low oncotic pressure. Chronic kidney disease also impairs VLDL clearance due to down-regulation of LPL and hepatic lipase.
    \item \textbf{Medications}: Many commonly prescribed drugs can raise triglycerides. These include oral estrogens, glucocorticoids, atypical antipsychotics, beta-blockers, thiazide diuretics, retinoids, cyclosporine, tacrolimus, and protease inhibitors.
\end{itemize}

\section*{Clinical Features}
In the majority of patients, mild-to-moderate hypertriglyceridemia (150--499 mg/dL) is entirely asymptomatic. The condition is most frequently discovered during routine lipid screening or during investigations for related metabolic disorders, such as type 2 diabetes or hypertension.

However, when triglyceride levels become severely elevated (typically exceeding 1000 mg/dL or 11.3 mmol/L), distinct clinical signs and symptoms can emerge due to the physical accumulation of lipids in tissues and the vascular system:
\begin{itemize}
    \item \textbf{Eruptive Xanthomas}: These are characteristic skin lesions that present as crops of small, 1--4 mm, yellowish-orange papules with an erythematous base. They typically appear suddenly on extensor surfaces such as the elbows, knees, buttocks, back, and hands. They consist of lipid-laden macrophages (foam cells) in the dermis and resolve slowly once triglyceride levels are lowered.
    \item \textbf{Lipemia Retinalis}: On fundoscopic examination of the eye, the retinal blood vessels may appear creamy, pinkish-white, or milky rather than their normal bright red color. This occurs when the concentration of light-scattering chylomicrons in the retinal capillaries is extremely high.
    \item \textbf{Hepatosplenomegaly}: Excess circulating chylomicrons are cleared by the reticuloendothelial system. The accumulation of these lipids in macrophages within the liver and spleen can cause enlargement of these organs, presenting clinically as abdominal fullness, discomfort, or upper quadrant tenderness.
    \textbf{Neurological Manifestations}: Extremely high triglyceride levels can cause cognitive symptoms, often described by patients as ``brain fog,'' memory impairment, lethargy, irritability, or difficulty concentrating. Some patients may also experience paresthesias (tingling or numbness) in the hands and feet.
    \item \textbf{Acute Pancreatitis}: This is the most serious and potentially life-threatening clinical presentation. It characterized by the sudden onset of severe, constant, boring pain in the epigastrium or upper abdomen that classically radiates to the back. The pain is usually accompanied by nausea, persistent vomiting, abdominal distension, fever, and tachycardia.
\end{itemize}

\section*{Differential Diagnosis}
In evaluating a patient with hypertriglyceridemia, the clinician must differentiate between primary genetic disorders, secondary acquired conditions, and other causes of acute abdominal symptoms. The differential diagnoses to consider include:
\begin{itemize}
    \item \textbf{Familial Chylomicronemia Syndrome (FCS)}: A rare monogenic disorder presenting in childhood or early adulthood with severe, refractory hypertriglyceridemia (>1000 mg/dL) and recurrent pancreatitis, showing poor response to standard lipid-lowering drugs but responding well to severe dietary fat restriction.
    \item \textbf{Multifactorial Chylomicronemia Syndrome (MCS)}: A polygenic condition presenting in adulthood, where genetic susceptibility is triggered by secondary factors such as poorly controlled diabetes, alcohol abuse, or medications.
    \item \textbf{Familial Combined Hyperlipidemia (FCHL)}: Differentiated by the concurrent elevation of both LDL cholesterol and triglycerides, along with elevated ApoB levels and a strong family history of premature cardiovascular disease.
    \item \textbf{Dysbetalipoproteinemia (Type III Hyperlipoproteinemia)}: Differentiated by elevated cholesterol and triglycerides in roughly equal proportions, the presence of xanthoma striatum palmare (yellowish discoloration of the palmar creases), and genetic confirmation of the APOE2/E2 genotype.
    \item \textbf{Secondary Hyperlipidemia}: Identifying whether the dyslipidemia is secondary to underlying conditions such as uncontrolled diabetes mellitus, severe hypothyroidism, nephrotic syndrome, or Cushing's disease.
    \item \textbf{Non-Lipidic Acute Pancreatitis}: If presenting with acute abdominal pain and elevated pancreatic enzymes, other causes of pancreatitis must be ruled out, including gallstones (cholelithiasis), alcohol-induced pancreatitis, hypercalcemia, drug-induced pancreatitis (e.g., azathioprine, valproate), and anatomical abnormalities of the pancreatic duct.
\end{itemize}

\section*{When to Seek Medical Care}
Fasting triglyceride levels should be monitored and discussed with a primary care physician during routine health examinations. Patients should schedule a medical appointment if they receive a lipid panel report indicating triglyceride levels above 150 mg/dL, or if they experience unexplained, recurrent mild abdominal discomfort, changes in vision, or the appearance of yellow skin bumps.

\textbf{Emergency Medical Situations}:
Severely elevated triglycerides can precipitate acute pancreatitis, which is a medical emergency. Immediate emergency medical services must be contacted if a patient experiences any of the following symptoms:
\begin{itemize}
    \item Sudden, severe, and unrelenting abdominal pain, particularly in the upper abdomen, that radiates through to the back.
    \item Inability to tolerate oral fluids due to persistent, severe vomiting.
    \item High fever associated with rapid breathing or a racing heart rate.
    \item Dizziness, confusion, cold or clammy skin, or loss of consciousness, which are signs of systemic shock.
\end{itemize}

For life-threatening symptoms, immediately call the local emergency services number:
\begin{itemize}
    \item \textbf{000} in many countries.
    \item \textbf{911} in North America (United States and Canada).
    \item \textbf{999} in the United Kingdom.
    \item \textbf{112} in the European Union and many other international regions.
\end{itemize}
Patients experiencing these severe symptoms should not attempt to drive themselves to the hospital but should wait for emergency medical technicians to arrive, as critical supportive therapy may need to begin immediately.

\section*{Diagnosis and Investigations}
The diagnosis of hypertriglyceridemia is established by measuring fasting plasma triglyceride levels. The diagnostic workup is directed toward confirming the elevation, categorizing its severity, assessing cardiovascular risk, and identifying any secondary causes.

The diagnostic process involves several steps and laboratory investigations:
\begin{itemize}
    \item \textbf{Fasting Lipid Profile}: While non-fasting samples are increasingly accepted for routine cardiovascular screening, a fasting lipid panel (requiring a 9--12 hour fast, during which only water is allowed) is the gold standard. Fasting is necessary because postprandial chylomicrons can temporarily double triglyceride levels, making accurate assessment difficult.
    \item \textbf{Diagnostic Thresholds}: Fasting triglyceride concentrations are classified as:
    \begin{itemize}
        \item Normal: <150 mg/dL (<1.7 mmol/L).
        \item Borderline High: 150--199 mg/dL (1.7--2.2 mmol/L).
        \item High: 200--499 mg/dL (2.3--5.6 mmol/L).
        \item Very High (Severe): >=500 mg/dL (>=5.6 mmol/L). Levels exceeding 880 mg/dL (10 mmol/L) are often classified as very severe, indicating a high risk of pancreatitis.
    \end{itemize}
    \item \textbf{Screening for Secondary Etiologies}: A comprehensive panel should include:
    \begin{itemize}
        \item Fasting plasma glucose and Hemoglobin A1c (HbA1c) to evaluate for prediabetes and diabetes.
        \item Thyroid-stimulating hormone (TSH) to screen for hypothyroidism.
        \item Serum creatinine, estimated glomerular filtration rate (eGFR), and urinalysis to check for protein excretion to rule out chronic kidney disease and nephrotic syndrome.
        \item Liver function tests (LFTs), including AST and ALT, to evaluate for metabolic dysfunction-associated steatotic liver disease (MASLD).
    \end{itemize}
    \item \textbf{Pancreatic Enzymes}: In cases of suspected pancreatitis, serum amylase and lipase must be measured immediately. Serum lipase is preferred because it remains elevated longer and has higher sensitivity. It is important to note that severe hypertriglyceridemia can interfere with some laboratory assays, leading to falsely normal or mildly elevated amylase levels; dilution of the serum sample is required in these cases.
    \item \textbf{Additional Lipoprotein Assays}: In patients with moderate hypertriglyceridemia, non-HDL cholesterol and apolipoprotein B (ApoB) should be measured. These markers provide a more accurate estimation of the total number of circulating atherogenic particles than LDL cholesterol alone.
    \item \textbf{Imaging and Genetics}: An abdominal ultrasound is useful to detect fatty liver or hepatosplenomegaly. Genetic testing may be performed to confirm primary genetic syndromes like FCS, particularly in young patients with severe elevations and recurrent pancreatitis in the absence of secondary factors.
\end{itemize}

\section*{Management}
The management strategy for hypertriglyceridemia is determined by the severity of the elevation. For patients with mild-to-moderate elevations, the primary goal is to reduce long-term cardiovascular risk. For patients with severe elevations, the immediate priority is to lower triglycerides below 500 mg/dL to prevent acute pancreatitis.

\textbf{Lifestyle Modifications}:
Lifestyle changes are the first-line therapy and remain a cornerstone of management:
\begin{itemize}
    \item \textbf{Dietary Changes}: For mild-to-moderate elevations, patients should reduce their intake of simple sugars, refined carbohydrates, and saturated fats, while adopting a Mediterranean-style diet rich in fiber, whole grains, vegetables, and polyunsaturated fats. For severe hypertriglyceridemia (>500 mg/dL), dietary fat must be severely restricted to less than 10\% to 15\% of total daily calories (approximately 15--20 grams of fat per day) to prevent the accumulation of dietary chylomicrons.
    \item \textbf{Alcohol Cessation}: Because alcohol is a potent stimulator of hepatic VLDL synthesis and inhibits LPL, complete abstinence is strongly recommended for all patients with hypertriglyceridemia, especially those with severe elevations.
    \item \textbf{Physical Activity}: Regular aerobic exercise, such as brisk walking, cycling, or swimming for at least 150 minutes per week, helps lower triglycerides by improving insulin sensitivity and enhancing LPL activity.
    \textbf{Weight Loss}: A reduction in body weight of 5\% to 10\% can lower triglyceride levels by 20\% to 30\% in individuals who are overweight or obese.
\end{itemize}

\textbf{Pharmacological Therapy}:
Medication is indicated when lifestyle changes are insufficient or when triglyceride levels are high enough to pose an immediate risk:
\begin{itemize}
    \item \textbf{Statins (HMG-CoA Reductase Inhibitors)}: Statins (e.g., atorvastatin, rosuvastatin) are the primary therapy for moderate hypertriglyceridemia when the clinical goal is cardiovascular risk reduction. While statins primarily target LDL cholesterol, they also lower triglycerides by 10\% to 30\% by increasing the clearance of VLDL remnants.
    \item \textbf{Fibrates}: Fibrates (e.g., fenofibrate, gemfibrozil) are peroxisome proliferator-activated receptor-alpha (PPAR-alpha) agonists that increase LPL activity, decrease hepatic VLDL secretion, and stimulate fatty acid oxidation. They typically lower triglycerides by 30\% to 50\%. Fibrates are the drug of choice for severe hypertriglyceridemia (>500 mg/dL) to prevent pancreatitis. Fenofibrate is preferred over gemfibrozil when combined with statins, as gemfibrozil significantly increases the risk of statin-induced myopathy.
    \item \textbf{Omega-3 Fatty Acids}: Prescription-strength omega-3 fatty acids (such as icosapent ethyl or mixed omega-3 acid ethyl esters) at a dose of 4 grams per day can reduce triglycerides by 20\% to 45\% by inhibiting hepatic VLDL synthesis. Highly purified icosapent ethyl has been shown in large clinical trials to reduce cardiovascular events in high-risk patients who have elevated triglycerides despite statin therapy.
    \item \textbf{Niacin (Nicotinic Acid)}: Niacin inhibits hepatic triglyceride synthesis and can lower triglycerides by 20\% to 50\%. However, its clinical use has declined significantly due to a lack of incremental cardiovascular benefit when added to statin therapy, as well as a high incidence of side effects, including cutaneous flushing, hyperglycemia, and hepatotoxicity.
    \item \textbf{Novel Therapies}: For patients with FCS, volanesorsen (an antisense oligonucleotide that inhibits apolipoprotein C-III) is approved in some jurisdictions to lower triglycerides through an LPL-independent pathway.
\end{itemize}

\textbf{Emergent Management of Hypertriglyceridemic Pancreatitis}:
In patients hospitalized with acute pancreatitis due to severe hypertriglyceridemia, specific interventions to rapidly clear triglycerides are required:
\begin{itemize}
    \item \textbf{Intravenous Insulin Infusion}: Regular insulin stimulates the expression and activity of LPL, promoting the rapid clearance of chylomicrons. This is administered under close monitoring, typically alongside a dextrose infusion to prevent hypoglycemia.
    \item \textbf{Therapeutic Plasma Exchange (Plasmapheresis)}: In severe cases, especially those with signs of organ failure or systemic inflammation, plasmapheresis is used to physically filter and remove chylomicrons and triglycerides from the blood, rapidly lowering levels to below 500 mg/dL.
\end{itemize}

\section*{Complications}
Untreated or poorly managed hypertriglyceridemia can lead to several serious health complications:
\begin{itemize}
    \item \textbf{Acute Pancreatitis}: This is the most urgent complication. The risk increases significantly when triglycerides exceed 500 mg/dL and becomes very high when levels exceed 1000 mg/dL. The proposed mechanism is that excess triglycerides are hydrolyzed by pancreatic lipase in the microcirculation of the pancreas, releasing large amounts of free fatty acids. These fatty acids exceed the binding capacity of albumin and exert a direct toxic effect on pancreatic acinar cells and capillary endothelium, leading to ischemia, inflammation, and necrosis.
    \item \textbf{Atherosclerotic Cardiovascular Disease (ASCVD)}: Moderate hypertriglyceridemia is an independent risk factor for coronary artery disease, myocardial infarction, ischemic stroke, and peripheral arterial disease. Triglyceride-rich remnant lipoproteins can penetrate the arterial wall, undergo oxidation, and promote foam cell formation and plaque development. Elevated triglycerides are also frequently associated with low HDL cholesterol and an abundance of small, dense LDL particles, which are highly atherogenic.
    \item \textbf{Metabolic Dysfunction-Associated Steatohepatitis (MASH)}: The liver response to excess triglycerides is the accumulation of lipid droplets within hepatocytes, leading to hepatic steatosis. In some patients, this lipotoxicity triggers chronic inflammation and cellular injury (steatohepatitis), which can progress to liver fibrosis, cirrhosis, and hepatic failure.
    \item \textbf{Chronic Pancreatitis and Endocrine/Exocrine Insufficiency}: Recurrent episodes of acute pancreatitis can lead to permanent structural damage to the pancreas. This can result in chronic pancreatitis, characterized by chronic pain, exocrine pancreatic insufficiency (causing malabsorption, steatorrhea, and weight loss), and endocrine pancreatic insufficiency (resulting in pancreatogenic or Type 3c diabetes).
\end{itemize}

\section*{Prognosis and Outcomes}
The prognosis for individuals with hypertriglyceridemia is variable and depends on the severity of the lipid elevation, the etiology, the presence of cardiovascular disease, and the patient's adherence to treatment.

For the majority of patients with mild-to-moderate hypertriglyceridemia, the prognosis is excellent if cardiovascular risk factors are managed. Dietary changes, regular exercise, and appropriate pharmacotherapy (such as statins or omega-3 fatty acids) are highly effective in lowering triglyceride levels, improving metabolic profile, and reducing the incidence of major adverse cardiovascular events.

For patients with severe or very severe hypertriglyceridemia, the prognosis is heavily influenced by the risk of acute pancreatitis. While the mortality rate of a single episode of hypertriglyceridemic pancreatitis ranges from 5\% to 10\%, it can rise significantly in patients who develop complications such as pancreatic necrosis or multi-organ failure. Patients with genetic disorders such as FCS face a chronic, lifelong risk of recurrent pancreatitis, which can lead to chronic pain and permanent pancreatic damage, severely impacting their quality of life. Long-term outcomes in these patients depend on their ability to adhere to a strict, low-fat diet.

\section*{Prevention and Self-Care}
Preventing the development of hypertriglyceridemia and managing existing elevations at home requires a combination of routine screening, healthy lifestyle habits, and adherence to medical therapy:
\begin{itemize}
    \item \textbf{Routine Lipid Screening}: Adults should undergo regular lipid panel screening starting at age 20, repeated every 4 to 6 years, or more frequently if risk factors such as obesity, diabetes, or a family history of lipid disorders are present.
    \item \textbf{Dietary Management}:
    \begin{itemize}
        \item Restrict the intake of simple sugars, sweetened beverages, and foods containing high-fructose corn syrup.
        \item Substitute refined grains with fiber-rich whole grains such as oats, brown rice, and quinoa.
        \item Emphasize healthy fats, such as monounsaturated and polyunsaturated fats found in olive oil, avocados, nuts, and fatty fish, while avoiding trans fats and limiting saturated fats.
        \item For individuals with severe hypertriglyceridemia, strictly limit total fat intake to less than 15 grams per day. Under medical supervision, medium-chain triglyceride (MCT) oil may be used for cooking, as MCTs are absorbed directly into the portal vein and do not require chylomicron transport.
    \end{itemize}
    \item \textbf{Exercise and Weight Control}: Engage in at least 150 minutes of moderate-intensity aerobic exercise per week. Maintain a healthy body weight, as weight loss is highly effective in improving insulin sensitivity and lowering triglyceride levels.
    \item \textbf{Alcohol Restriction}: Limit alcohol consumption or avoid it entirely, particularly if triglyceride levels are elevated, as alcohol is a potent trigger for VLDL production.
    \item \textbf{Glycemic Management}: If diagnosed with diabetes, maintain strict blood glucose control, as hyperglycemia directly impairs LPL activity and raises triglyceride levels.
    \item \textbf{Medication Adherence}: Take all prescribed lipid-lowering medications exactly as directed. Do not discontinue fibrates, statins, or omega-3 supplements without consulting a doctor, even if repeat lipid tests show normal levels.
\end{itemize}

\section*{Sources and Further Reading}
For further information regarding triglyceride metabolism, cardiovascular health, and lipid disorders, refer to the following authoritative organizations:
\begin{itemize}
    \item American Heart Association (AHA): Offers comprehensive guidelines and educational resources on lipids and cardiovascular disease prevention. Website: \texttt{https://www.heart.org}
    \item National Heart, Lung, and Blood Institute (NHLBI): A division of the U.S. National Institutes of Health providing clinical guidelines and health information on cholesterol and triglycerides. Website: \texttt{https://www.nhlbi.nih.gov}
    \item European Society of Cardiology (ESC): Publishes clinical practice guidelines for the management of dyslipidemias and prevention of cardiovascular disease. Website: \texttt{https://www.escardio.org}
    \item Mayo Clinic: Provides patient-centered resources explaining the causes, diagnosis, and treatment options for high triglycerides. Website: \texttt{https://www.mayoclinic.org}
    \item National Lipid Association (NLA): A professional society focused on the study and management of lipid disorders, offering materials for both patients and clinicians. Website: \texttt{https://www.lipid.org}
    \item The Endocrine Society: A global professional organization providing clinical guidelines on lipid disorders and metabolic syndrome. Website: \texttt{https://www.endocrine.org}
    \item Heart Foundation Australia: An independent charity providing practical dietary and lifestyle advice for cardiovascular health. Website: \texttt{https://www.heartfoundation.org.au}
\end{itemize}